\section*{Chapter \ref{ch:b2:angiosperm-reproduction} --- Sexual Reproduction of Flowering Plants}

\begin{solution}{exo:b2:angiosperm-reproduction:1}
Sepals (calyx), petals (corolla), \omterm{def:b2:angiosperm-reproduction:flower}{stamens} (androecium), \omterm{def:b2:angiosperm-reproduction:flower}{carpels}
(gynoecium). \omterm{def:b2:angiosperm-reproduction:flower}{Stamen}: filament and anther. \omterm{def:b2:angiosperm-reproduction:flower}{Carpel}: \omterm{def:b2:angiosperm-reproduction:flower}{stigma}, style,
\omterm{def:b2:angiosperm-reproduction:flower}{ovary} with \omterm{def:b2:plant-life-cycles:heterospory}{ovules}. All these are \omterm{def:b2:meiosis-heredity:meiosis}{diploid} \omterm{def:b2:plant-life-cycles:cycles}{sporophyte}; only the \omterm{prop:b2:angiosperm-reproduction:pollen}{pollen
grain} (inside the anther) and the \omterm{prop:b2:angiosperm-reproduction:embryosac}{embryo sac} (inside the \omterm{def:b2:plant-life-cycles:heterospory}{ovule}) are
\omterm{def:b2:meiosis-heredity:meiosis}{haploid} \omterm{def:b2:plant-life-cycles:cycles}{gametophyte}.
\end{solution}

\begin{solution}{exo:b2:angiosperm-reproduction:2}
\omterm{def:b2:plant-life-cycles:heterospory}{Microspore} $n$; vegetative cell $n$; sperm $n$; egg $n$; polar
nucleus $n$; zygote $2n$; endosperm $3n$; \omterm{def:b2:angiosperm-reproduction:seed}{seed coat} $2n$ (maternal
integuments); \omterm{def:b2:angiosperm-reproduction:seed}{fruit} wall $2n$ (maternal \omterm{def:b2:angiosperm-reproduction:flower}{ovary}).
\end{solution}

\begin{solution}{exo:b2:angiosperm-reproduction:3}
Wind (grasses, oak, birch, hazel, plantain): small green \omterm{def:b2:angiosperm-reproduction:flower}{flowers}, no
petals or scent or \omterm{def:b2:angiosperm-reproduction:pollination}{nectar}, anthers dangling on long filaments,
feathery \omterm{def:b2:angiosperm-reproduction:flower}{stigmas}, vast amounts of smooth dry \omterm{def:b2:plant-life-cycles:heterospory}{pollen}, flowering before
the leaves. Bee (clover, sage, foxglove, lavender, apple): coloured
petals often with ultraviolet guides, landing platform, scent,
\omterm{def:b2:angiosperm-reproduction:pollination}{nectar}, sticky sculptured \omterm{def:b2:plant-life-cycles:heterospory}{pollen} in moderate amounts, flowering in
daylight.
\end{solution}

\begin{solution}{exo:b2:angiosperm-reproduction:4}
The two sperm cells of one \omterm{thm:b2:angiosperm-reproduction:tube}{pollen tube} fuse, one with the egg
(zygote, $2n$, the embryo) and one with the central cell (endosperm
nucleus, $3n$, the food store). The endosperm begins to grow only
after fertilisation, so the plant provisions only \omterm{def:b2:plant-life-cycles:heterospory}{ovules} that hold
an embryo, unlike the gymnosperm, which builds its food store before.
\end{solution}

\begin{solution}{exo:b2:angiosperm-reproduction:5}
$25/1.2 = \qty{21}{h}$. Silks emerging on days 5, 6 and 7 meet \omterm{def:b2:plant-life-cycles:heterospory}{pollen};
those of days 8 to 12 do not: $3/8$ of the silks are pollinated, and
the ear is five-eighths barren.
\end{solution}

\begin{solution}{exo:b2:angiosperm-reproduction:6}
$S_1S_2\times S_1S_2$: 0. $S_1S_2\times S_1S_3$: $\tfrac12$ (only $S_3$
\omterm{def:b2:plant-life-cycles:heterospory}{pollen} grows). $S_1S_2\times S_3S_4$: all. Offspring of the second
cross: eggs $S_1$ or $S_2$, sperm $S_3$: $S_1S_3$ and $S_2S_3$, half
each.
\end{solution}

\begin{solution}{exo:b2:angiosperm-reproduction:7}
A plant carries 2 of the $n$ \omterm{def:b2:meiosis-heredity:vocabulary}{alleles} and rejects \omterm{def:b2:plant-life-cycles:heterospory}{pollen} bearing
either: compatible fraction $(n - 2)/n$. A new \omterm{def:b2:meiosis-heredity:vocabulary}{allele} is rejected by
no pistil, so its \omterm{def:b2:plant-life-cycles:heterospory}{pollen} sires more \omterm{def:b2:angiosperm-reproduction:seed}{seeds} than average and it spreads
until it is as common as the rest; the same advantage protects every
rare \omterm{def:b2:meiosis-heredity:vocabulary}{allele} from loss. Selection therefore accumulates \omterm{def:b2:meiosis-heredity:vocabulary}{alleles}, and
natural populations of self-incompatible species carry dozens.
\end{solution}

\begin{solution}{exo:b2:angiosperm-reproduction:8}
The two \omterm{prop:b2:angiosperm-reproduction:embryosac}{polar nuclei} are copies of one \omterm{def:b2:plant-life-cycles:heterospory}{megaspore}, so the central
cell is $YY$ or $yy$ (half each); the sperm is $Y$ or $y$ (half
each): endosperm $YYY$, $YYy$, $Yyy$, $yyy$ at $\tfrac14$ each; three
quarters yellow. Pollinated by $yy$: $YYy$ and $yyy$, half each ---
half the kernels yellow.
\end{solution}

\begin{solution}{exo:b2:angiosperm-reproduction:9}
Barley grains are cut in half; the embryo-bearing halves digest
their starch, the embryo-less halves do not, unless gibberellin is
added, when they do. (a) The embryo is the source of the signal, not
of the enzyme. (b) Gibberellin is the signal, sufficient by itself.
(c) The amylase is made by the aleurone layer of the endosperm,
which responds to the hormone. The embryo-less half separates signal
from response: without it one could not tell whether the embryo
digests the starch itself.
\end{solution}

\begin{solution}{exo:b2:angiosperm-reproduction:10}
Wind: $10^{6}\times\qty{1}{\micro g} = \qty{1}{g}$ per \omterm{def:b2:plant-life-cycles:heterospory}{ovule}. Animal:
\qty{1}{mg} of \omterm{def:b2:plant-life-cycles:heterospory}{pollen} plus \qty{0.5}{mg} of sugar $= \qty{1.5}{mg}$
per \omterm{def:b2:plant-life-cycles:heterospory}{ovule} --- about seven hundred times cheaper. Wind \omterm{def:b2:angiosperm-reproduction:pollination}{pollination}
needs no partner: it works in early spring before insects fly, in
cold, windy and open places, at night and in rain, in dense stands
of one species, and it cannot be cheated by \omterm{def:b2:angiosperm-reproduction:pollination}{nectar} thieves or
abandoned by a pollinator that goes extinct.
\end{solution}

\begin{solution}{exo:b2:angiosperm-reproduction:11}
\omterm{def:b2:angiosperm-reproduction:seed}{Seeds} falling at the foot of the parent land in its shade, among its
roots, and where its specific \omterm{def:b2:angiosperm-reproduction:seed}{seed} predators and pathogens
concentrate; almost all die. A bird carries a \omterm{def:b2:angiosperm-reproduction:seed}{seed} kilometres
(\qty{12}{km} in twenty minutes) and drops it, with fertiliser, in a
hedge or a clearing; even if only one \omterm{def:b2:angiosperm-reproduction:seed}{seed} in a hundred lands in a
good site, that is more than all the \omterm{def:b2:angiosperm-reproduction:seed}{seeds} under the tree, and the
population spreads at the speed of birds, not of falling \omterm{def:b2:angiosperm-reproduction:seed}{fruit}. The
flesh is the price of the ticket.
\end{solution}

\begin{solution}{exo:b2:angiosperm-reproduction:12}
Enclosure makes the \omterm{def:b2:plant-life-cycles:heterospory}{pollen} grow through maternal tissue, which
allows the pistil to test it (self-incompatibility, rejection of
foreign species), to let many tubes compete so the fastest sires the
\omterm{def:b2:angiosperm-reproduction:seed}{seed}, to protect \omterm{def:b2:plant-life-cycles:heterospory}{ovules} from drying and from herbivores, and to turn
the \omterm{def:b2:angiosperm-reproduction:flower}{ovary} into a \omterm{def:b2:angiosperm-reproduction:seed}{fruit} that recruits animals for dispersal; \omterm{prop:b2:angiosperm-reproduction:double}{double
fertilisation} then provisions only fertilised \omterm{def:b2:plant-life-cycles:heterospory}{ovules}. The costs: a
\omterm{prop:b2:angiosperm-reproduction:pollen}{pollen grain} must reach a \omterm{def:b2:angiosperm-reproduction:flower}{stigma} rather than the \omterm{def:b2:plant-life-cycles:heterospory}{ovule} itself, so a
style, a tube and its guidance are needed, and the \omterm{def:b2:angiosperm-reproduction:seed}{fruit} is a heavy
investment. The balance favoured angiosperms so strongly that they
went from nothing to nine tenths of plant species in a hundred
million years.
\end{solution}

\begin{solution}{pb:b2:angiosperm-reproduction:1}
\textbf{1.} On A ($S_1S_2$): \omterm{def:b2:plant-life-cycles:heterospory}{pollen} A $f = 0$; B ($S_1$, $S_3$) $f =
\tfrac12$; C ($S_3$, $S_4$) $f = 1$.
\textbf{2.} On B ($S_1S_3$): A $\tfrac12$, B 0, C $\tfrac12$. On C
($S_3S_4$): A 1, B $\tfrac12$, C 0.
\textbf{3.} 50 compatible grains; $0.1\times 50 = 5$ \omterm{def:b2:plant-life-cycles:heterospory}{ovules}.
\textbf{4.} From C: 100 compatible, $0.1\times 100 = 10$, i.e.\ all
ten \omterm{def:b2:plant-life-cycles:heterospory}{ovules}; from A: none.
\textbf{5.} \omterm{def:b2:angiosperm-reproduction:seed}{Fruits} from C visits (10 \omterm{def:b2:angiosperm-reproduction:seed}{seeds}) are kept; from B (5 \omterm{def:b2:angiosperm-reproduction:seed}{seeds})
just kept; from A dropped. An orchard of A alone bears nothing.
\textbf{6.} Half the \omterm{def:b2:angiosperm-reproduction:flower}{flowers} set \omterm{def:b2:angiosperm-reproduction:seed}{fruit} (those visited from C): 2500
\omterm{def:b2:angiosperm-reproduction:seed}{fruits} per tree, ten times the 250 needed --- a full crop, and the
tree will drop the surplus.
\textbf{7.} It sheds compatible \omterm{def:b2:plant-life-cycles:heterospory}{pollen} over the whole flowering season
of every variety; but if it shares both $S$ \omterm{def:b2:meiosis-heredity:vocabulary}{alleles} with a variety
it pollinates nothing on it.
\textbf{8.} $8\times 10^{7}/(7\times 86400) = 132$ grains per square
metre per second.
\textbf{9.} $132/0.25 = 530$ grains per cubic metre.
\textbf{10.} $132\times 10^{-4} = 0.013$ per second: one grain every
\qty{76}{s}.
\textbf{11.} $0.013\times 86400 \approx 1100$ grains a day; the first
tube to reach the \omterm{def:b2:plant-life-cycles:heterospory}{ovule} fertilises it, and the \omterm{def:b2:plant-life-cycles:heterospory}{ovule} is then closed
to the rest.
\textbf{12.} \qty{20}{h} after \omterm{def:b2:angiosperm-reproduction:pollination}{pollination}.
\textbf{13.} \omterm{def:b2:plant-life-cycles:heterospory}{Pollen} on days 0--7, silks from day 5: only days 5, 6 and
7 overlap --- $3/7$ of the shedding period, and silks emerging after
day 7 get nothing; the ear is largely barren. Drought at silking is
the classic cause of a failed maize crop.
\textbf{14.} Embryo $Yy$; endosperm $YYy$.
\textbf{15.} Yellow, since $Y$ is \omterm{def:b2:meiosis-heredity:vocabulary}{dominant}: the \omterm{def:b2:plant-life-cycles:heterospory}{pollen} parent's trait
shows in the \omterm{def:b2:angiosperm-reproduction:seed}{seed} on the mother plant --- xenia.
\textbf{16.} Central cell $YY$ or $yy$: endosperm $YYy$ (half) and
$yyy$ (half).
\textbf{17.} $YYy$: 20 units; $Yyy$: 10 units; a $yyY$ with the $Y$
from the father is the same \omterm{def:b2:meiosis-heredity:vocabulary}{genotype} $Yyy$, 10 units. Two paternal
doses cannot exist: one sperm fertilises the central cell, and the
endosperm is always two maternal genomes to one paternal.
\textbf{18.} It is the starch store; $0.3/0.33 = \qty{91}{\%}$ of the
kernel.
\textbf{19.} The store is built only after fertilisation, so no
resources go into \omterm{def:b2:plant-life-cycles:heterospory}{ovules} that carry no embryo; the gymnosperm's
\omterm{def:b2:plant-life-cycles:cycles}{gametophyte} is built beforehand and wasted if \omterm{def:b2:angiosperm-reproduction:pollination}{pollination} fails.
\textbf{20.} $0.9\times 600 = 540$ kernels per ear and per plant;
$4320$ per square metre; $4320\times 0.3 = \qty{1.3}{kg}$ of endosperm
per square metre.
\textbf{21.} \qty{13}{t/ha}.
\textbf{22.} Grain carbon $0.5\times 1.2 = \qty{0.6}{kg/m^2}$;
endosperm carbon $0.45\times 1.3 = \qty{0.58}{kg/m^2}$: consistent
(the small remainder is embryo and coat).
\textbf{23.} $8\times 10^{7}/4320 \approx \num{18500}$ grains per
kernel.
\textbf{24.} Its own \omterm{def:b2:plant-life-cycles:heterospory}{pollen} cloud is thin and blown away, and its
tassel sheds mostly before its own silks emerge, so many silks are
never pollinated and the ear has gaps.
\textbf{25.} A block: half the \omterm{def:b2:angiosperm-reproduction:flower}{flowers} set \omterm{def:b2:angiosperm-reproduction:seed}{fruit}, 2500 per tree, a
full crop. Field: 4320 kernels per square metre, \qty{13}{t/ha}.
\end{solution}
